% from ega4-10-fr.tex
\oldpage[IV-3]{109}
On sait en effet que $\dim(U)\leq\dim(X)$ ($\mathbf 0$, 14.1.4)~; en outre,
comme $X$ est un préschéma de Jacobson, $U$ contient un point fermé de toute
composante irréductible de $X$, donc $\dim(U)=\dim(X)$ en vertu de (10.6.1) et
($\mathbf 0$, 14.1.2.1).

\begin{proposition}[10.6.3]
\label{IV.10.6.3-fr}
Supposons que $X$ vérifie les conditions $1^\circ$, $2^\circ$ et $3^\circ$ de
(10.6.1), et soit $Y$ une partie fermée de $X$. Alors, pour tout $x\in Y$, et
tout voisinage ouvert $U$ de $x$ dans $X$ ne rencontrant pas les composantes
irréductibles de $Y$ qui ne contiennent pas $x$, on a
\begin{equation}
\label{IV.10.6.3.1-fr}
\tag{10.6.3.1}
\dim_x(Y)=\dim(U\cap Y)=\dim(\overline{\{x\}})
 +\operatorname{codim}(\overline{\{x\}},Y)=\sup_i\dim(Y_i)
\end{equation}
où $Y_i$ ($1\leq i\leq m$) sont les composantes irréductibles de $Y$ contenant
$x$.
\end{proposition}

En considérant un sous-préschéma fermé de $X$ ayant $Y$ pour espace sous-jacent,
on peut se borner, en vertu de (10.6.1, (i)), au cas où $Y=X$. En vertu du choix
de $U$, $U$ est réunion des $U\cap X_i$, donc
$\dim(U)=\sup_i\dim(U\cap X_i)$~; mais d'après (10.6.2) appliqué à un
sous-préschéma fermé de $X$ ayant $X_i$ pour espace sous-jacent, on a (compte
tenu de (10.6.1, (i))) $\dim(U\cap X_i)=\dim(X_i)$~; cela démontre que le second
et le quatrième terme de (10.6.3.1) sont égaux. Comme $U\cap X_i$ est
biéquidimensionnel (10.6.1, (i)) (puisque l'immersion $U\cap X_i\to X_i$ est de
type fini (\textbf{I}, 6.3.5)), on a
\begin{equation}
\label{IV.10.6.3.2-fr}
\tag{10.6.3.2}
\dim(X_i)=\dim(U\cap X_i)=\dim(U\cap\overline{\{x\}})
 +\operatorname{codim}(U\cap\overline{\{x\}},U\cap X_i)
\end{equation}
($\mathbf 0$, 14.3.5). D'après (10.6.2) appliqué à un sous-préschéma fermé de
$X$ ayant $\overline{\{x\}}$ pour espace sous-jacent,
$\dim(U\cap\overline{\{x\}})=\dim(\overline{\{x\}})$~; comme on a aussi, pour les
mêmes raisons,
\begin{equation}
\label{IV.10.6.3.3-fr}
\tag{10.6.3.3}
\dim(X_i)=\dim(\overline{\{x\}})
 +\operatorname{codim}(\overline{\{x\}},X_i)
\end{equation}
on obtient
\[
  \dim(U)=\dim(\overline{\{x\}})
  +\sup_i\operatorname{codim}(\overline{\{x\}},X_i)
  =\dim(\overline{\{x\}})+\operatorname{codim}(\overline{\{x\}},X)
\]
par définition de la codimension ($\mathbf 0$, 14.2.1). Ceci montre que $\dim(U)$
est indépendante du voisinage ouvert $U$ de $x$ vérifiant les conditions de
l'énoncé, donc est égal à $\dim_x(X)$, par ($\mathbf 0$, 14.1.4.1).

\begin{corollary}[10.6.4]
\label{IV.10.6.4-fr}
Sous les hypothèses de (10.6.3), soient $\mathcal F$ un $\mathcal O_X$-Module cohérent,
$Y$ son support. Pour tout $x\in Y$, on a
\begin{equation}
\label{IV.10.6.4.1-fr}
\tag{10.6.4.1}
\dim(\overline{\{x\}})+\dim(\mathcal F_x)=\dim_xY.
\end{equation}
\end{corollary}

En effet, cela résulte de (10.6.3.1) et de la formule
$\dim(\mathcal F_x)=\operatorname{codim}(\overline{\{x\}},Y)$ (5.1.12.2).

\subsection{Exemples et contre-exemples}
\label{subsection:IV.10.7-fr}

\begin{env}[10.7.1]
\label{IV.10.7.1-fr}
Soit $S$ un préschéma localement noethérien de dimension $\leq1$ et supposons
que $S$ soit un \emph{préschéma de Jacobson}~; lorsque $S$ est noethérien, il
revient au même de dire que les composantes irréductibles de $S$ de dimension
$1$ sont infinies, car tout $x\in S$ qui n'est pas fermé est le point générique
d'une telle composante (10.4.5 et 10.5.4). Alors $S$ vérifie aussi les conditions
$2^\circ$ et $3^\circ$ de (10.6.1)~: en effet, tout anneau local $\mathcal O_{S,s}$
est de dimension $0$ ou $1$, et par suite est universellement caténaire (7.2.9)~;
d'autre part, une composante irréductible $S'$ de $S$ est, soit réduite

